\scanpage{183}
\mspage{57}
% archive sheet number: 182

La \unsure{dernière|deuxième} assertion de 5.7. résulte aussitôt du fait qu'une limite inductive filtrante de monos dans $\mathcal M$ est mono. (condition e) p. 46)

On peut préciser la prop. 5.6. obtenue ainsi :

\uline{Proposition 5.7.} \deleted{\unsure{Soient}} Soit $\mathcal M$ une catégorie satisfaisant les conditions \unsure{a) b)} c) d) (p. 44,45\unsure{,|-} et e) (p. 46). Soit $M$ une partie de $\mathrm{Ob}\,\mathcal M$, stable par isomorphie, \deleted{les conditions \unsure{suivantes} \illegible{1}} \inserted{$\unsure{\pi|\Pi}$ un cardinal tel} que les objets de $\mathrm{Filt}^{\Pi}$ soient $\pi$-accessibles. Les conditions suivantes sont équivalentes :

a) $M$ est $\pi$-accessible, et stable par $\varinjlim$ filtrantes

b) $M$ est $\pi$ accessible, et $\inserted{M\cap\mathrm{Filt}^{\Pi}}$ \deleted{\unsure{est}} est stable par $\underline{\lim}$ filtrantes \deleted{\unsure{de}} cardinal $\leq\pi$.

% slanted
\margin{\deleted{Condition de \unsure{$\pi$-compacité} qui peut se vérifier \unsure{maintenant} en \unsure{5.1)} p. 50, cf en bas..} \unsure{$\to$}}

c) Il existe une partie $M^\pi$ de $\mathrm{Filt}^{\Pi}$, stable par $\underline{\lim}$ filtrantes de cardinal $\leq\pi$, \deleted{\unsure{et pour}} et telle que pour tout $X\in\mathrm{Ob}\,\mathcal M$, \deleted{\unsure{toute} partie $J$ cofinale de $\mathrm{Ss}_\pi(X)$ (stable par bornes supérieures de parties filtrantes de cardinal $\leq\pi$)} on ait \deleted{\unsure{équivalence}} $X\in M$ sss l'ens $J_M$ des $Z_\alpha\in J$ tels que $Z_\alpha\in M^\pi$ \deleted{\unsure{soit}} est cofinal dans $J$.

(De plus, sous la condition c) on a

